
\section{二阶线性常微分方程（齐次）}
我们现在转向本章的主要内容——二阶线性常微分方程。这类方程极为重要，
因为它们出现在量子力学、电磁理论以及物理学其他领域中求解偏微分方程的常用方法里。
与一阶线性常微分方程不同，二阶线性常微分方程通常没有普适的封闭解，
一般需要采用幂级数形式的解法。在讨论级数解法的一般方法之前，我们首先考察微分方程中奇点的概念。

奇点

微分方程的奇点概念对我们有两方面的重要意义：
（1）它有助于对微分方程进行分类，并识别那些可以转化为常见形式的方程
（本小节后面会讨论）；（2）它关系到能否找到级数解的可行性。

% 解的可行性，这正是 Fuchs 定理（稍后将讨论）的主题。

当一个线性齐次二阶常微分方程写成如下形式时：

$$
y''+p(x) y '+q(x) y=0
$$

若 $p(x)$ 和 $q(x)$ 在 $x_0$ 处有限，则称 $x_0$ 为该方程的常点（ordinary point）。如果 $p(x)$ 或 $q(x)$ 在 $x \rightarrow x_0$ 时发散，则 $x_0$ 称为奇点（singular point）。奇点进一步分为正则奇点（regular）和非常奇点（irregular，也称为本性奇点 essential singularity）：
\begin{itemize}
    \item 若 $x_0$ 为奇点，但 $\left(x-x_0\right) p(x)$ 和 $\left(x-x_0\right)^2 q(x)$ 在 $x_0$ 处仍有限，则 $x_0$ 为正则奇点。
    \item 若 $p(x)$ 发散速度快于 $1 /\left(x-x_0\right)$，使得 $\left(x-x_0\right) p(x)$ 在 $x \rightarrow x_0$ 时趋于无穷，或 $q(x)$ 发散快于 $1 /\left(x-x_0\right)^2$，使得 $\left(x-x_0\right)^2 q(x)$ 在 $x \rightarrow x_0$ 时趋于无穷，则 $x_0$ 为非常奇点。
\end{itemize}

这些定义适用于所有有限的 $x_0$。
\begin{note}
    注意这里的$\infty$为常奇点的情况具体
    解释我们并没有给出, 它的解决需要用$z=1/x$的代换研究$z=0$时的情形, 具体讨论这里略去.
\end{note}

% 若要分析 $x \rightarrow \infty$ 处的行为，可令 $x=1 / z$，代入微分方程，
% 并考察 $z \rightarrow 0$ 的极限。此时原本关于 $y(x)$ 的方程将转化为关于 $w(z)$ 的方程，
% 其中 $w(z)=y\left(z^{-1}\right)$。对导数进行如下变换，

% $$
% y''+p(x) y '+q(x) y=0
% $$
% $x_0$ 时, $p\left(x_0\right), q\left(x_0\right)$ 有限,称 $x_0$为ODE的常点
% $p\left(x_0\right), q\left(x_0\right) \rightarrow \infty$称 $x_0$为ODE的奇点

常奇点(或正则奇点) $$\left\{\begin{array}{l}
    \left(x-x_0\right) p\left(x_0\right) \text{有限}
    \\
    \left(x-x_0\right)^2 q(x) \text{有限}
    \\ p\left(x_0\right), q(x) \rightarrow \infty .\end{array}
    \right.$$
其他情况为非常奇点.


下面给出几种重要的ODE的常奇点
\begin{enumerate}
    \item 超几何(hypergeometic):
    $$
    x(x-1) y''+[(1+a+b) x+c] y'+a b y= 0
    $$
    常奇点 $0,1, \infty$\\
    \item 勒让德 (Legendre):
    $$
    \left(1-x^2\right) y''-2 x y'+l(l+1) y=0
    $$
    常奇点. $-1,1, \infty$ .\\
    \item  切比雪夫. (Chebyshev):
    $$
    \left(1-x^2\right) y''-x y'+n^2 y=0 
    $$
    常奇点 $-1,1, \infty$.\\

    \item  合流超几何 (confluent hypergeometic):
    $$
    x y''+(c-x) y'-a y=0,
    $$
    \\
    \item 拉盖尔(Laguerre)
    $$
    x y''+(1-x) y'+a y=0,
    $$
    常奇点 $0$.
    \item 贝塞尔(Bessel)
    $$
    x^2 y'' + xy' + (x^2 - n^2) y = 0
    $$
    常奇点$0$.
\end{enumerate}


\subsection{几类特殊的二阶方程}
\subsubsection{ $y'' = f(x)$ 型}
此类方程只需要积分两次就可以得到通解
\begin{equation}
    y = \int F(x) dx + C_1 x + C_2
\end{equation}
其中$F(x) = \int f(x) dx$, $C_1, C_2$是积分常数.
\subsubsection{ $y'' = f(x, y')$ 型}
利用换元$u = y'$, 则有$u' = y''$, 原方程化为一阶方程
$$ u' = f(x,u)$$
该解为$u = \varphi(x, C_1)$, 原方程解为
\begin{equation}
    y = \int \varphi(x, C_1) dx  + C_2.
\end{equation}
% \subsubsection{ $y'' = f(y, y')$ 型}


\subsection{常点邻域上的级数解}

考虑线性（经典）振子方程
$$
\frac{d^2 y}{d x^2}+\omega^2 y=0
$$

我们已经用其他方法解过这个方程，得到的解为 $y=\sin \omega x$ 和 $\cos \omega x$。

我们尝试设

$$
\begin{aligned}
y(x) & =x^s\left(a_0+a_1 x+a_2 x^2+a_3 x^3+\cdots\right) \\
& =\sum_{j=0}^{\infty} a_j x^{s+j}, \quad a_0 \neq 0
\end{aligned}
$$

其中指数 $s$ 以及所有系数 $a_j$ 都尚未确定。注意 $s$ 不必为整数。对 $y(x)$ 两次求导，得到

$$
\begin{aligned}
\frac{d y}{d x} & =\sum_{j=0}^{\infty} a_j(s+j) x^{s+j-1}, \\
\frac{d^2 y}{d x^2} & =\sum_{j=0}^{\infty} a_j(s+j)(s+j-1) x^{s+j-2} .
\end{aligned}
$$

将上述表达式代入方程，有

$$
\sum_{j=0}^{\infty} a_j(s+j)(s+j-1) x^{s+j-2}+\omega^2 \sum_{j=0}^{\infty} a_j x^{s+j}=0 .
$$




% $\left|x-x_0\right|<R$ 上.$x_0$为常点.

% $$
% y(x)=\sum_{k=0} a_k\left(x-x_0\right)^k, a_1, a_2, a_k \cdots \text { 待定. }
% $$
% 这里以谐振子为例
% $$
% \text { 令 } \begin{aligned}
% y(x) & =x^s\left(a_0+a_1 x+a_2 x^2+\cdots\right) \\
% & =\sum_{j=0}^{\infty} a_j x^{s+j}, a_0 \neq 0 .
% \end{aligned}
% $$

% \begin{note}
% 这里课堂上采用了直接令$s=0$的讲法.
% \end{note}
% $$
% \begin{aligned}
% & \frac{d y}{d x}=\sum_{j=0}^{\infty} a_j(j+s) x^{s+j-1} \\
% & \frac{d^2 y}{d x^2}=\sum_{j=0}^{\infty} a_j(j+s)(j+s-1) x^{s+j-2}
% \end{aligned}
% $$

% 代入原方程.
% $$
% \sum_{j=0}^{\infty} a_j(s+j)(s+j-1) x^{s+j-2}+\omega^2 \sum_{j=0}^{\infty} a_j x^{s+j=0} \text {. }
% $$
由于幂级数的性质, 各个系数必须为零.
% $$
% \text { 任意阶 } x^{s+j} \text { 都满足 }
% $$
% ,系数必为零.
对于最低幂次项有
$$
 a_0(s)(s-1)=0 
$$
这个方程来自于最低次幂$x$的系数，被称为判别方程（indicial equation）。
判别方程及其根对我们的分析至关重要。在本例中，它清楚地告诉我们$s=0$或$s=1$，
因此我们的级数解必须以$x^0$或$x^1$项开始。

进一步考察, 我们发现下一个最低次幂$x$，即$x^{s-1}$，
也只在一项中出现（即第一组求和中$j=1$时）。令$x^{s-1}$的系数为零，有
$$
a_1(s+1)s=0 .
$$
这表明如果$s=1$，则必须有$a_1=0$。而如果$s=0$，该方程对系数没有任何限制。

$$
\begin{aligned}
% & x^{s-2} a_0(s)(s-1)=0 \quad\left(a_0 \neq 0\right) \\
& x^{s-1}: a_1(s+1)(s)=0 \\
& x^s: a_2(s+2)(s+1)+a_0 w^2=0
\end{aligned}
$$

若 $s=0$, $y \sim  a_0+a_1 x+\cdots$
若 $s=1, y \sim a_0 x+a_1 x^2+\cdots$
当 $s=1$时, $a_1=0$.
当 $s=0$ 时, $a_1$ 任意, 因此我们取$a_1=0$.
$$
\begin{gathered}
\Rightarrow a_{j+2}(s+j+2)(s+j+1)+w^2 a_j=0 \\
\Rightarrow a_{j+2}=\frac{-w^2}{(s+j+2)(s+j+1)} a_j
\end{gathered}
$$

递推关系.
若$s=0$ 
$$
\begin{aligned}
a_{j+2}=\frac{-\omega^2}{(j+2)(j+1)} a_j \\
& a_2=-\frac{w^2}{2 !} a_0 \\
& a_4=-\frac{a_2}{3^3 4}=+\frac{w^4}{4 !} a_0 \\
& a_6=-\frac{a_0^{3.4}}{5.6}=-\frac{w 6}{6 !} a_0 \\
&
\end{aligned}
$$

$$
a_{2 n}=(-1)^n \frac{w^{2 n}}{2 n !} a_0
$$

$$
\begin{aligned}
y & \left.=a_0\left[1-\frac{(\omega x)^2}{2 !}+\frac{(\omega x}{4}\right)^4-\frac{\omega x x^6}{6 !}+\omega\right] \\
& =a_0 \cos \omega x .
\end{aligned}
$$
若 $s=1$ 
$$
\begin{aligned}
& a_{j+2}=-a_j \frac{\omega^2}{(j+3) (j+2)} \\
& \Rightarrow a_2=-a_0 \frac{\omega^2}{2 \cdot 3}=\frac{\omega^2}{3 !} a_0 \\
& a_4=-a_2 \frac{\omega^2}{5 \cdot 4}=\frac{\omega^4}{5 !} a_0 \\
& a_6=-a_4 \frac{\omega^2}{6\cdot 7}=\frac{w^6}{7 !} a_0 \\
\end{aligned}
$$

$$
\begin{aligned}
& \Rightarrow \quad a_0 x\left[1-\frac{(\omega x)^2}{3 !}+\frac{(\omega x)^{4}}{5 !}-\frac{(\omega x)^6}{7 !} \cdots\right] \\
& =\frac{a_0}{\omega}\left[(\omega x)-\frac{(\omega x)^3}{3 !}+\frac{(\omega x)^5}{5 !} \cdots\right] \\
& =\frac{a_0}{\omega} \sin \omega x
\end{aligned}
$$
此方法称为Frobenius方法,上面关于$x_0$上展开的,一般的
我们可以在 $x_0$处展开
$$
y(x)=\sum_{j=0}^{\infty} a_j\left(x-x_0\right)^{s+j}, a_0 \neq 0 .
$$


\subsection{正则奇点邻域上的级数求解}
对线性振子的处理似乎过于简单。我们将幂级数（见式(7.28)）代入微分方程（见式(7.27)），很容易就得到了两个线性无关的解。

为了了解可能出现的其他情况，我们尝试求解贝塞尔方程：

$$
x^2 y^{\prime \prime}+x y^{\prime}+\left(x^2-n^2\right) y=0 .
$$

同样，假设解的形式为

$$
y(x)=\sum_{j=0}^{\infty} a_j x^{s+j},
$$

对其求导并代入方程（7.40），得到

$$
\begin{aligned}
& \sum_{j=0}^{\infty} a_j(s+j)(s+j-1) x^{s+j}+\sum_{j=0}^{\infty} a_j(s+j) x^{s+j} \\
& \quad+\sum_{j=0}^{\infty} a_j x^{s+j+2}-\sum_{j=0}^{\infty} a_j^2 x^{s+j}=0
\end{aligned}
$$

令 $j=0$，得到 $x^s$ 的系数，即左边最低次幂的系数：

$$
a_0\left[s(s-1)+s-n^2\right]=0
$$

且 $a_0 \neq 0$（定义如此）。因此，式（7.42）给出判别方程

$$
s^2-n^2=0
$$

解为 $s= \pm n$。
还需要考察 $x^{s+1}$ 的系数，得到

$$
a_1\left[(s+1) s+s+1-n^2\right]=0,
$$

即

$$
a_1(s+1-n)(s+1+n)=0 .
$$

对于 $s= \pm n$，$s+1-n$ 和 $s+1+n$ 都不为零，因此必须有 $a_1=0$。
继续考察 $x^{s+j}$ 的系数（取 $s=n$），它出现在第一、二、四项中 $a_j$ 的系数，在第三项中则是 $a_{j-2}$。令 $x^{s+j}$ 的总系数为零，得到

$$
a_j\left[(n+j)(n+j-1)+(n+j)-n^2\right]+a_{j-2}=0 .
$$

将 $j$ 替换为 $j+2$，对 $j \geq 0$ 可写为

$$
a_{j+2}=-a_j \frac{1}{(j+2)(2 n+j+2)},
$$

这就是所需的递推关系。不断应用该递推关系可得

$$
\begin{aligned}
& a_2=-a_0 \frac{1}{2(2 n+2)}=-\frac{a_0 n!}{2^2 1!(n+1)!} \\
& a_4=-a_2 \frac{1}{4(2 n+4)}=\frac{a_0 n!}{2^4 2!(n+2)!} \\
& a_6=-a_4 \frac{1}{6(2 n+6)}=-\frac{a_0 n!}{2^6 3!(n+3)!}, \quad \text { 以此类推 }
\end{aligned}
$$

一般地，

$$
a_{2 p}=(-1)^p \frac{a_0 n!}{2^{2 p} p!(n+p)!} .
$$

将这些系数代入假设的级数解，有

$$
y(x)=a_0 x^n\left[1-\frac{n!x^2}{2^2 1!(n+1)!}+\frac{n!x^4}{2^4 2!(n+2)!}-\cdots\right] .
$$

用求和式表示为

$$
\begin{aligned}
y(x) & =a_0 \sum_{j=0}^{\infty}(-1)^j \frac{n!x^{n+2 j}}{2^{2 j} j!(n+j)!} \\
& =a_0 2^n n!\sum_{j=0}^{\infty}(-1)^j \frac{1}{j!(n+j)!}\left(\frac{x}{2}\right)^{n+2 j} .
\end{aligned}
$$

在第14章中，最终的求和式（取 $a_0=1 / 2^n n!$）被定义为贝塞尔函数 $J_n(x)$：

$$
J_n(x)=\sum_{j=0}^{\infty}(-1)^j \frac{1}{j!(n+j)!}\left(\frac{x}{2}\right)^{n+2 j} .
$$

注意，这个解 $J_n(x)$ 具有偶或奇对称性，正如从贝塞尔方程的形式可以预期的那样。
当 $s=-n$ 且 $n$ 不是整数时，可以得到第二个不同的级数解，记为 $J_{-n}(x)$。但当 $-n$ 是负整数时，会出现问题。此时系数的递推关系仍由式（7.45）给出，只是 $2n$ 替换为 $-2n$。当 $j+2=2n$ 或 $j=2(n-1)$ 时，$a_{j+2}$ 的分母为零，Frobenius 方法无法得到与假设 $x^{-n}$ 开头的级数解相容的解。

通过代入无穷级数，我们为线性振子方程得到了两个解，为贝塞尔方程得到了一个（若 $n$ 不是整数则有两个）解。对于“是否总能这样做？该方法是否总是有效？”的问题，答案是“不，总不能这样做。级数解法并不总是有效。”


% 级数解有限个负幂项称为正则奇点 $p(x): \sum_{k=-1}^{\infty}$
% $q(x) : \sum_{k=-2}^{\infty}$

% 有判定方程,
% $$
% s(s-1)+s p_{-1}+q_{-2}=0
% $$
% 其中$s_1, s_2$ 为两根($s_1> s_2$).

% $$
% \begin{aligned}
% y_1(x) & =\sum_{k=0}^{\infty} a_k\left(x-x_0\right)^{s_1+k} \cdot \\
% y_2(x) & =\sum_{k=0}^{\infty} b_k\left(x-x_0\right)^{s_2+k} . \\
% \text{或} y_2(x) & =A y_1(x) \ln \left(x-x_0\right) \\
% & +\sum_{k=0}^{\infty} b_k\left(x-x_0\right)^{s_2+k}
% \end{aligned}
% $$


$\nu$阶贝塞尔方程的级数求解参考书上.




Frobenius 方法并不总是能给出两个解，当特征根的差为整数时,级数解法会给出两个线性相关的, 因此需要其他方法.
一般会给出一个级数解，找出第二个解的方式可以通过如下方法获得.

这里引入朗斯基行列式(Wronskian)
$$
\Delta(x)=\left|\begin{array}{ll}
y_1(x) & y_2(x) \\
y_1'(x) & y_2'(x)
\end{array}\right|= y_1(x) y_2'(x) - y_2(x) y_1'(x)
$$

由于$y_1, y_2$满足微分方程,
$$
\begin{aligned}
& \left(y_1 y_2''-y_1'' 
y_2\right)+p\left(y_1 y_2'-y_1' y_2\right) = 0\\
& \frac{d \Delta}{d x} - p \Delta(x)=0
\end{aligned}
$$

$$
\Delta(x)=\Delta_0 e^{-\int p(x) d x}
$$

$$
\begin{aligned}
& \frac{d}{d x}\left(\frac{y_{2}}{y_1}\right)=\frac{y_{1} y_{2}'-y_{1}' y_{2}}{y_{1}^2}=\frac{\Delta(x)}{y_{1}^2} \\
& y_{2}=y_{1} \int \frac{\Delta(x)}{\left(y_{1}(x)\right)^2} d x
\end{aligned}
$$

